%% Dans ce fichier, les �ventuels accents sont cod�s
%% en ISO-8859-15 (aka latin9)
%% La ligne suivante donne trois exemples 
%% e accent aigu : �, e accent grave : �, e-dans-l'o : �
%%%%%ENTREES%%%%
\begin{tabarg}
\multicolumn{3}{l}{\titretab {Entr\'ees}}\\
\hline
$X$ & Vecteur, \texttt{float} & S�rie de donn�es � filtrer (unit� quelconque).\\
$Y$ & Vecteur, \texttt{U\_Int\_8} & vecteur sur lequel seront appliqu�s les crit�res (unit� quelconde)\\
$Crit1_{min}$ & Coeff., \texttt{U\_Int\_8} & valeur mini de la premi�re gamme valide\\
$Crit1_{max}$ & Coeff., \texttt{U\_Int\_8} & valeur maxi de la premi�re gamme valide\\
$Crit2_{min}$ & Coeff., \texttt{U\_Int\_8} & valeur mini de la deuxi�me gamme valide\\
$Crit2_{max}$ & Coeff., \texttt{U\_Int\_8} & valeur maxi de la deuxi�me gamme valide\\
\end{tabarg}

\medskip
%%%%%SORTIES%%%%
\begin{tabarg}
\multicolumn{3}{l}{\titretab {Sorties}}\\
\hline
$X_{sel}$ & Vecteur, \texttt{float} & S�rie de donn�es filtr�es (m�me unit� que $X$)\\
\end{tabarg}
